\documentclass[11pt]{article}

% Standard, non-exotic packages only.
\usepackage[utf8]{inputenc}
\usepackage[T1]{fontenc}
\usepackage{amsmath, amssymb, amsthm}
\usepackage{mathtools}
\usepackage{booktabs}
\usepackage{array}
\usepackage[margin=1in]{geometry}
\usepackage{hyperref}
\hypersetup{colorlinks=true, linkcolor=blue, citecolor=blue, urlcolor=blue}

% Theorem environments.
\newtheorem{theorem}{Theorem}
\newtheorem{lemma}{Lemma}
\newtheorem{corollary}{Corollary}
\newtheorem{proposition}{Proposition}
\theoremstyle{definition}
\newtheorem{definition}{Definition}
\newtheorem{remark}{Remark}

% Shorthands.
\newcommand{\R}{\mathbb{R}}
\newcommand{\Q}{\mathbb{Q}}
\newcommand{\Z}{\mathbb{Z}}
\newcommand{\chrom}{\chi}
\DeclareMathOperator{\Aut}{Aut}

\title{Forcing-sterility of the realizable unit-distance lineage\\
and a codegree obstruction to $\chi \ge 6$}

% TODO before submission: confirm author name spelling and affiliation.
\author{Owen Kent\thanks{Independent researcher.
  Correspondence: \texttt{owenpkent@gmail.com}.}}
\date{June 2026}

\begin{document}
\maketitle

\begin{abstract}
The chromatic number of the plane, $\chrom(\R^2)$, is known only to lie in
$[5,7]$. The lower bound $\chrom(\R^2) \ge 5$ rests on finite unit-distance
graphs (UDGs): de Grey (2018) exhibited the first such graph with $\chrom \ge 5$,
and the Polymath16 project reduced the order to roughly $510$ vertices. A natural
route to $\chrom(\R^2) \ge 6$ would locate, somewhere in this realizable lineage
of $\chrom = 5$ UDGs, a non-adjacent vertex pair forced to distinct colors in
every proper $5$-coloring (a \emph{clamp}); realizing such a pair lifts the
chromatic number to $6$. This note records a negative structural result: the
realizable lineage contains no such forcing, and provably so. We recall an
elementary lemma (the \emph{Essential-Pair Lemma}, essentially known) which
implies that a vertex-critical $5$-chromatic graph cannot host any forced
non-adjacent pair. We confirm the
consequence by an exhaustive computational census: all $1{,}955{,}948$
non-adjacent pairs across the nine $\chrom = 5$ lineage graphs are free, with zero
forced pairs of either kind. Since the lineage is vertex-critical, this is a
one-line corollary of the Lemma, and the deeper point is conceptual: the lineage
was produced by SAT-driven minimization, which drives graphs toward
vertex-criticality, and criticality deletes exactly the vertices that could host
forcing. The lineage is forcing-sterile \emph{by construction}. We then frame the
search for the missing $\chrom \ge 6$ object through a codegree obstruction: every
planar UDG is $K_{2,3}$-free, so its edge count is bounded by
$m \le n\bigl(1 + \sqrt{8n-7}\bigr)/4$; combined with the Kostochka-Yancey lower
bound for $6$-critical graphs and the vertex Folkman number
$F_v(2,2,2,2,2;4) = 16$, this excludes every small dense Folkman-type host. The
missing object must lie in the class of graphs that are simultaneously $K_4$-free
and $K_{2,3}$-free, of order at least roughly $26$. We claim no new bound on
$\chrom(\R^2)$.
\end{abstract}

\section{Introduction}
\label{sec:intro}

The Hadwiger-Nelson problem asks for the chromatic number $\chrom(\R^2)$ of the
\emph{unit-distance graph} on the Euclidean plane: vertices are the points of
$\R^2$, and two points are adjacent iff they are at distance exactly $1$. The
classical bounds are
\[
  5 \le \chrom(\R^2) \le 7,
\]
and have resisted improvement for three quarters of a century. The upper bound
$\chrom(\R^2) \le 7$ is the Isbell hexagonal coloring (tile the plane with regular
hexagons of diameter slightly less than $1$ and color in a repeating
$7$-pattern). The lower bound $\chrom(\R^2) \ge 5$ is recent and computational:
de Grey~\cite{deGrey2018} constructed a finite unit-distance graph on $1581$
vertices with $\chrom \ge 5$, verified by a SAT proof, and the Polymath16
project~\cite{Polymath16} subsequently reduced the order of such graphs to
roughly $510$ vertices. We refer to this family of explicitly realized $\chrom=5$
UDGs as the \emph{realizable lineage}. For background on the problem and its
history see Soifer~\cite{Soifer2009}.

The natural finitary route to $\chrom(\R^2) \ge 6$ is to find a finite UDG with
$\chrom \ge 6$. One concrete mechanism for manufacturing such a graph out of a
$\chrom=5$ graph is a \emph{clamp}.

\begin{definition}[Clamp]
\label{def:clamp}
Let $G$ be a $5$-chromatic graph and $(s,t)$ a non-adjacent pair of vertices. We
call $(s,t)$ a \emph{clamp} if every proper $5$-coloring of $G$ assigns
$c(s) \ne c(t)$.
\end{definition}

If a clamp $(s,t)$ in a UDG can be \emph{realized}, that is, if the geometry
permits placing a vertex at unit distance from both $s$ and $t$ while preserving
the unit-distance structure (so that an edge effectively appears across the
forced-different pair), the chromatic number is lifted to $6$. The motivating
question for this note is therefore:
\begin{quote}
\emph{Does the realizable lineage of $\chrom = 5$ UDGs contain a clamp (or any
other forced non-adjacent pair) that could be exploited to reach
$\chrom \ge 6$?}
\end{quote}
The answer of this note is \textbf{no, and provably so}. We make this precise
through one elementary lemma (Section~\ref{sec:lemma}), confirm it with an
exhaustive computational census (Section~\ref{sec:census}), explain it
structurally by way of vertex-criticality (Section~\ref{sec:criticality}), and
then reframe the search for the genuinely missing object through a codegree
obstruction (Section~\ref{sec:codegree}). Section~\ref{sec:discussion}
synthesizes the picture and states the target object honestly.

\paragraph{Scope and honesty.} This is a negative, explanatory note. We move no
bound on $\chrom(\R^2)$. The contribution is threefold: (a) the observation that
the realizable lineage is forcing-sterile \emph{by construction}, obtained by
applying a known elementary lemma to its vertex-criticality; (b) an exhaustive
computational census confirming it across the entire lineage; and
(c) a codegree framing that locates the missing $\chrom \ge 6$ object outside the
reach of both the realizable lineage and the small dense graphs one would
naively manufacture. None of the three ingredients is individually deep; the
synthesis is the point.

\section{Preliminaries}
\label{sec:prelim}

All graphs are finite, simple, and undirected. For a graph $G$ we write $V(G)$
for its vertex set, $n = |V(G)|$, $m = |E(G)|$, and $\chrom(G)$ for the chromatic
number. A graph is \emph{$5$-chromatic} if $\chrom(G) = 5$. A proper
$k$-coloring is a map $c : V(G) \to \{1, \dots, k\}$ with $c(u) \ne c(v)$ whenever
$uv \in E(G)$.

\begin{definition}[Vertex-critical]
A graph $G$ is \emph{vertex-critical} (with respect to its chromatic number) if
$\chrom(G - v) < \chrom(G)$ for every $v \in V(G)$. A $5$-chromatic graph is
vertex-critical iff $\chrom(G - v) \le 4$ for every vertex $v$.
\end{definition}

\begin{definition}[Forcing of a non-adjacent pair]
\label{def:forcing}
Let $G$ be a $5$-chromatic graph and let $(s,t)$ be a pair of distinct,
\emph{non-adjacent} vertices. We classify $(s,t)$ as follows.
\begin{itemize}
  \item \textsc{Forced-same}: every proper $5$-coloring of $G$ has $c(s) = c(t)$.
  \item \textsc{Forced-different}: every proper $5$-coloring of $G$ has
        $c(s) \ne c(t)$.
  \item \textsc{Forced}: either \textsc{forced-same} or \textsc{forced-different}
        holds.
  \item \textsc{Free}: neither holds, i.e. there exist a proper $5$-coloring with
        $c(s) = c(t)$ and a (possibly different) proper $5$-coloring with
        $c(s) \ne c(t)$.
\end{itemize}
\end{definition}

A clamp (Definition~\ref{def:clamp}) is exactly a forced-different non-adjacent
pair. A forced-same pair is no weaker as a route to $\chrom \ge 6$: adding a
single unit edge across a forced-same pair produces a forced-different pattern
elsewhere (a splice argument), so forced-same and forced-different are
interconvertible as $\chrom \ge 6$ ingredients. For the purposes of this note the
relevant dichotomy is simply \emph{forced} versus \emph{free}: a single forced
pair anywhere in the lineage would be a candidate $\chrom \ge 6$ ingredient.

\paragraph{Deciding forcing by SAT.} Each forcing question reduces to two
$k$-colorability decisions, both decidable by a standard $k$-coloring SAT
encoding.
\begin{itemize}
  \item \textsc{Forced-different} holds for $(s,t)$ iff the graph $G / \{s,t\}$,
        obtained by \emph{identifying} (merging) $s$ and $t$, is not
        $5$-colorable. Merging forces $c(s) = c(t)$; if even that constrained
        instance is $5$-UNSAT, then no proper $5$-coloring has $c(s) = c(t)$, so
        every proper $5$-coloring has $c(s) \ne c(t)$.
  \item \textsc{Forced-same} holds for $(s,t)$ iff the graph $G + st$, obtained by
        \emph{adding} the edge $st$, is not $5$-colorable. Adding the edge forbids
        $c(s) = c(t)$; if that instance is $5$-UNSAT, then no proper $5$-coloring
        has $c(s) \ne c(t)$, so every proper $5$-coloring has $c(s) = c(t)$.
  \item Otherwise $(s,t)$ is \textsc{free}.
\end{itemize}
Both tests are decisive: a $5$-UNSAT verdict from a complete solver is a proof,
and a satisfying assignment is an explicit certificate of freeness on that side.

\paragraph{Control objects.} Throughout we keep in mind the structural controls
of the program. The unit-distance graph on $\Q^2$ has chromatic number $2$
(Woodall~\cite{Woodall1973}); any purely combinatorial argument for
$\chrom(\R^2) \ge 5$ must use the topology of $\R$ in a way that fails on $\Q^2$.
The codegree obstruction of Section~\ref{sec:codegree} is genuinely Euclidean:
it uses two-circle rigidity, which fails under the sup-norm (where unit
``circles'' are squares and can share a segment). The Essential-Pair Lemma of
Section~\ref{sec:lemma}, by contrast, is pure graph theory and claims no bound
anywhere.

\section{The Essential-Pair Lemma}
\label{sec:lemma}

The following lemma is elementary and essentially known. Its corollary, that a
vertex-critical $k$-chromatic graph admits no forced non-adjacent pair, appears
as Theorem~3.17 (``there are no implicit relations in critical graphs'') of
Martin~\cite{Martin2009}, where forced-different and forced-same pairs are termed
\emph{implicit-edges} and \emph{implicit-identities}; the single-deletion form we
state is a one-line sharpening of his Theorems~3.2 and~3.17. The result is also
implicit in standard critical-graph theory (Dirac and Gallai; see Jensen and
Toft~\cite{JensenToft1995}). We state and prove it for self-containedness and
make no priority claim; the proof uses only the definition of a proper coloring.

\begin{lemma}[Essential-Pair Lemma]
\label{lem:essential}
Let $G$ be a $5$-chromatic graph and let $(s,t)$ be a non-adjacent pair that is
\emph{forced} (forced-same or forced-different). Then
\[
  \chrom(G - s) \ge 5 \quad\text{and}\quad \chrom(G - t) \ge 5.
\]
\end{lemma}

\begin{proof}
We prove $\chrom(G - s) \ge 5$; the statement for $t$ follows by symmetry.

Suppose toward a contradiction that $\chrom(G - s) \le 4$. Fix a proper
$4$-coloring of $G - s$ using colors $\{1,2,3,4\}$. We extend it to $G$ in two
different ways.

\emph{The ``different'' pattern.} Give $s$ the new color $5$. Every neighbor of
$s$ lies in $G - s$ and so is colored with a color in $\{1,2,3,4\}$; hence no
neighbor of $s$ has color $5$, and the extension is a proper $5$-coloring of $G$.
In it, $c(s) = 5$. Since $t \ne s$ and $t$ is colored with a color in
$\{1,2,3,4\}$, we have $c(s) \ne c(t)$. Thus the forced-different pattern is
realized.

\emph{The ``same'' pattern.} Starting again from the same proper $4$-coloring of
$G - s$, set $c(s) = 5$ as before, and additionally \emph{recolor} $t$ to $5$.
Because $s$ and $t$ are \textbf{non-adjacent}, every neighbor of $t$ also lies in
$G - s$ (a neighbor of $t$ cannot be $s$), and so is colored in $\{1,2,3,4\}$;
recoloring $t$ to $5$ therefore creates no monochromatic edge at $t$. The vertex
$s$ likewise has all neighbors in $\{1,2,3,4\}$, so no edge at $s$ is
monochromatic either. The result is a proper $5$-coloring of $G$ with
$c(s) = 5 = c(t)$. Thus the forced-same pattern is realized.

Both the same-colored and the differently-colored patterns occur in proper
$5$-colorings of $G$. Hence $(s,t)$ is \textsc{free}
(Definition~\ref{def:forcing}), contradicting the hypothesis that it is forced.
Therefore $\chrom(G - s) \ge 5$, and by symmetry $\chrom(G - t) \ge 5$.
\end{proof}

\begin{corollary}
\label{cor:critical-free}
A vertex-critical $5$-chromatic graph contains \emph{no} forced non-adjacent
pair: every non-adjacent pair is free.
\end{corollary}

\begin{proof}
Let $G$ be vertex-critical and $5$-chromatic, so $\chrom(G - v) \le 4$ for every
vertex $v$. If some non-adjacent pair $(s,t)$ were forced, then
Lemma~\ref{lem:essential} would give $\chrom(G - s) \ge 5$, contradicting
$\chrom(G - s) \le 4$. Hence no non-adjacent pair is forced.
\end{proof}

The content of Corollary~\ref{cor:critical-free} is intuitive once stated: a
forced pair $(s,t)$ requires that the rest of the graph be ``heavy'' enough at
both $s$ and $t$ to constrain their colors even without an edge between them, and
that heaviness is precisely what survives the deletion of $s$ (or $t$). If the
graph is already minimal in the sense that every vertex deletion drops the
chromatic number, no such residual heaviness exists.

\section{The exhaustive census}
\label{sec:census}

We now report the computational confirmation. The realizable lineage of
$\chrom = 5$ UDGs studied here consists of nine graphs, of orders
\[
  510,\; 517,\; 529,\; 553,\; 610,\; 633,\; 803,\; 826,\; 874,
\]
descending from the de Grey / Polymath16 constructions. (An earlier triage table
listed three further graphs, on $199$, $403$, and $721$ vertices, which on
re-examination have $\chrom = 4$; they are excluded here, and their exclusion
does not affect the census headline.)

\paragraph{Method.} For every non-adjacent pair $(s,t)$ in each of the nine
graphs we ran \emph{both} decisive SAT tests of Section~\ref{sec:prelim}:
\begin{itemize}
  \item forced-different is decided by \emph{merging} $s$ and $t$ (identifying the
        two vertices) and testing whether the merged graph is $5$-UNSAT;
  \item forced-same is decided by \emph{adding} the edge $st$ and testing whether
        the augmented graph is $5$-UNSAT;
  \item if neither test returns UNSAT, the pair is free.
\end{itemize}
A per-solve conflict budget was used and tuned so that the number of
indeterminate solves was driven to exactly zero; every pair received a decisive
classification.

\paragraph{Result.} All $1{,}955{,}948$ non-adjacent pairs across the nine graphs
were classified, and \textbf{every single one is free}: zero forced-same, zero
forced-different, zero indeterminate. Table~\ref{tab:census} reports the
per-graph counts.

\begin{table}[ht]
\centering
\caption{Exhaustive forcing census over the realizable $\chrom = 5$ lineage. For
each graph, every non-adjacent pair was classified by two decisive SAT tests
(merge for forced-different, add-edge for forced-same). Every pair is free.}
\label{tab:census}
\begin{tabular}{r r r r r}
\toprule
Order $n$ & Non-adjacent pairs & Forced-same & Forced-diff. & Free \\
\midrule
510 & 127{,}291 & 0 & 0 & all \\
517 & 130{,}807 & 0 & 0 & all \\
529 & 136{,}986 & 0 & 0 & all \\
553 & 149{,}906 & 0 & 0 & all \\
610 & 182{,}745 & 0 & 0 & all \\
633 & 196{,}862 & 0 & 0 & all \\
803 & 317{,}859 & 0 & 0 & all \\
826 & 336{,}452 & 0 & 0 & all \\
874 & 377{,}040 & 0 & 0 & all \\
\midrule
\textbf{Total} & \textbf{1{,}955{,}948} & \textbf{0} & \textbf{0} &
  \textbf{all} \\
\bottomrule
\end{tabular}
\end{table}

% Per-graph non-adjacent-pair counts are C(n,2) - m, computed directly from the
% graph edge files; they sum to 1,955,948 over the nine chi=5 graphs. (The
% figure 2,309,264 from an earlier run counted twelve graphs, including the
% three chi=4 graphs S199/L403/T721 now excluded; that broader total is not the
% nine-graph headline.)

Equivalently: in the realizable lineage, every non-adjacent pair admits both a
same-colored and a differently-colored proper $5$-coloring. All color forcing is
\emph{adjacency-local}. A single non-free pair would have been a candidate
$\chrom \ge 6$ ingredient (a forced-different pair is a clamp directly; a
forced-same pair yields one by a splice). None exists.

\section{The criticality explanation}
\label{sec:criticality}

The census of Section~\ref{sec:census} is not a coincidence, and it is not even
a separate fact. A criticality scan over the same nine graphs (deciding
$4$-colorability of $G - v$ for every vertex $v$, totaling on the order of
$5{,}000$ solves) found that \emph{all nine lineage graphs are vertex-critical}:
deleting any single vertex yields a $4$-colorable graph, with zero redundant
vertices and zero indeterminate verdicts.

\begin{theorem}
\label{thm:sterile}
Every graph in the realizable $\chrom = 5$ lineage is vertex-critical, and
therefore (by Corollary~\ref{cor:critical-free}) contains no forced non-adjacent
pair. Consequently the exhaustive census of Section~\ref{sec:census} is a
one-line corollary of the Essential-Pair Lemma.
\end{theorem}

\begin{proof}
Vertex-criticality is the computational output of the scan. Given it,
Corollary~\ref{cor:critical-free} applies to each of the nine graphs and yields
that no non-adjacent pair is forced, which is exactly the census conclusion.
\end{proof}

\paragraph{The conceptual heart.} Theorem~\ref{thm:sterile} explains the census
mechanistically, and the explanation is the central point of this note. The
realizable lineage was \emph{produced by SAT-driven minimization}: starting from
a large $\chrom = 5$ UDG, one repeatedly deletes vertices and edges while
preserving $\chrom = 5$, driving the graph toward a minimal core. Minimization of
this kind drives a graph toward vertex-criticality. But by the Essential-Pair
Lemma, criticality deletes \emph{exactly} the vertices that could host a forced
pair: a forced pair $(s,t)$ certifies that $s$ and $t$ are chromatically
\emph{redundant} (their deletion keeps $\chrom \ge 5$), and a vertex-critical
graph has no redundant vertices. The production process therefore removes every
clamp-capable endpoint before any forcing search is ever run.

\begin{quote}
\emph{The realizable lineage is forcing-sterile by construction.} The missing
$\chrom \ge 6$ ingredient cannot be hiding in it, not because no one searched
carefully, but because the very procedure that produced the lineage removed the
only vertices that could carry the ingredient.
\end{quote}

This has a concrete methodological consequence. A clamp (or any forced pair) can
only live at chromatically redundant vertices, so a candidate-generation
procedure aimed at $\chrom \ge 6$ must \emph{preserve} redundancy: it must
\textbf{not} minimize toward criticality. Equivalently, any forcing sweep can be
pre-filtered by an $O(n)$ criticality scan (keep only redundant vertices) rather
than testing all $O(n^2)$ pairs, since forcing can only occur at redundant
endpoints. We note one caveat carried from prior work: redundancy is necessary
but not sufficient. There exist non-minimal $5$-chromatic graphs (for instance a
particular $C_6$-closure on $1155$ vertices) that are forcing-free despite having
redundant vertices.

\section{A codegree obstruction to small dense hosts}
\label{sec:codegree}

If the missing $\chrom \ge 6$ object is not in the realizable lineage
(Sections~\ref{sec:census}-\ref{sec:criticality}), the alternative is to
manufacture an abstract $\chrom = 6$ graph and realize it as a UDG. The natural
supply is $K_4$-free $6$-chromatic graphs (one wants $K_4$-freeness because unit
distances make $K_4$ geometrically rigid and hard to realize freely). These exist
in abundance and can be made small. The obstruction recorded here is that they
are \emph{too dense} to be planar UDGs.

\paragraph{The codegree bound.} In the Euclidean plane, two distinct unit circles
meet in at most $2$ points. Hence any two vertices of a planar UDG have at most
$2$ common neighbors: \textbf{every planar UDG is $K_{2,3}$-free}. (This is the
classical fact underlying the Spencer-Szemer\'edi-Trotter $O(n^{4/3})$ bound on
the number of unit distances among $n$ points~\cite{SpencerSzemerediTrotter1984}.)

The $K_{2,3}$-free condition bounds the edge count. Counting paths of length two
through each vertex,
\[
  \sum_{v} \binom{d_v}{2} \;=\; \#\{\text{paths of length } 2\}
  \;\le\; 2\binom{n}{2},
\]
because each unordered pair of vertices is the pair of endpoints of at most $2$
such paths ($K_{2,3}$-freeness caps the common-neighbor count at $2$). By
convexity, $\sum_v \binom{d_v}{2} \ge n \binom{2m/n}{2}$, and solving the
resulting quadratic in $m$ gives
\begin{equation}
  m \;\le\; \frac{n\bigl(1 + \sqrt{8n - 7}\bigr)}{4}.
  \label{eq:codeg-ceiling}
\end{equation}
We call the right-hand side the \emph{codegree ceiling}.

\paragraph{Manufactured hosts overshoot it.} Empirically, the smallest reachable
$K_4$-free $6$-chromatic (Folkman-type) hosts are dense, with edge-to-vertex
ratio $m/n \approx 4.2$, far above the ceiling
of~\eqref{eq:codeg-ceiling}. Concretely:
\begin{itemize}
  \item an $18$-vertex $K_4$-free $6$-critical host has $m = 76$, against a
        ceiling of $\lfloor 57.17 \rfloor = 57$;
  \item the classical Mycielskian $M^3(C_5)$ on $47$ vertices has $434$ vertex
        pairs of common-neighbor count $\ge 3$, i.e. $434$ violations of the
        $K_{2,3}$-free condition, with maximum codegree $11$.
\end{itemize}
Such graphs cannot be planar UDGs; the violation is structural (it persists for
edge-critical cores, which have no slack to remove), not an artifact of any
particular construction heuristic.

\paragraph{The violations are intrinsic to the critical core.} One might hope that
the $K_{2,3}$ violations of a manufactured host are removable ``padding'' that a
careful sparsification could shed while preserving $\chrom = 6$. A top-down probe
indicates otherwise. Starting from $M^3(C_5)$ (the $47$-vertex $\chrom = 6$
Mycielskian above, with $434$ violating pairs) and repeatedly deleting edges while
holding $\chrom = 6$ as a hard constraint---re-establishing $6$-chromaticity by a
compensating forced-same edge whenever a deletion would drop it---drives the graph
down to $m = 131$ edges, essentially the Kostochka-Yancey $6$-critical floor
$2.8 \cdot 47 - 1.8 = 129.8$. Yet the codegree excess does not fall to zero: the
resulting edge-critical, $K_4$-free, $\chrom = 6$ graph still has $166$ vertex pairs
of common-neighbor count $\ge 3$. Sparsifying all the way to the critical floor
removes none of the $K_{2,3}$ violations; they live entirely in the $6$-critical
core. For this construction, then, the chromatic obstruction and the codegree
obstruction are the \emph{same} substructure: one cannot shed the codegree
violations without destroying the very $6$-chromaticity they accompany. (This is a
heuristic observation about the $M^3(C_5)$-derived family, obtained by local search,
not a theorem about all $6$-critical graphs.)

\paragraph{A rigorous core-size floor.} The $K_4$-free assumption can be dropped
for a clean lower bound on the order of the $6$-critical core. The
Kostochka-Yancey theorem~\cite{KostochkaYancey2014} states that every
$6$-critical graph has
\[
  m \;\ge\; 2.8\,n - 1.8.
\]
Comparing this floor with the codegree ceiling~\eqref{eq:codeg-ceiling}, the
floor \emph{exceeds} the ceiling for $n \le 12$ (and they cross at $n = 13$).
Hence:

\begin{proposition}
\label{prop:core13}
Any $6$-chromatic unit-distance graph in the plane contains a $6$-critical core
on at least $13$ vertices.
\end{proposition}

\begin{proof}
A $6$-chromatic graph contains a $6$-critical subgraph $H$ (delete vertices and
edges while $\chrom = 6$). As a subgraph of a planar UDG, $H$ is $K_{2,3}$-free,
so $|E(H)| \le |V(H)|\bigl(1 + \sqrt{8|V(H)| - 7}\bigr)/4$. As a $6$-critical
graph, $|E(H)| \ge 2.8\,|V(H)| - 1.8$. For $|V(H)| \le 12$ the lower bound
exceeds the upper bound, a contradiction, so $|V(H)| \ge 13$.
\end{proof}

\paragraph{The Folkman floor and the empty window.} The minimum order of a
$K_4$-free $6$-chromatic graph is the vertex Folkman number
$F_v(2,2,2,2,2;4)$. The literature settles this at
\[
  F_v(2,2,2,2,2;4) = 16
\]
(see Xu-Liang-Radziszowski~\cite{XuLiangRadziszowski2018}). Its only $16$-vertex
witnesses are the two Ramsey $(4,4)$-graphs (Paley-type), which have $m \approx 60$.
Both violate the codegree ceiling, since~\eqref{eq:codeg-ceiling} gives exactly
$48$ at $n = 16$ and $60 > 48$. Therefore the UDG-necessary class, that is, the
class of graphs that are simultaneously $K_4$-free \emph{and} $K_{2,3}$-free, has
no member at the Folkman floor. Because the ceiling~\eqref{eq:codeg-ceiling} only
reaches average degree $7.5$ (the density of the Ramsey $(4,4)$-graphs) around
$n \approx 26$, even Folkman-floor density excludes the UDG-necessary class below
roughly $n \approx 26$.

\paragraph{Novelty.} We stress that the ingredients here are not new. The
$K_{2,3}$-free fact is classical (it is the combinatorial heart of the unit-distance
incidence bound~\cite{SpencerSzemerediTrotter1984}), the Kostochka-Yancey bound
is a published theorem, and the Folkman number is a literature value. The
contribution of this section is the \emph{synthesis}: pinning the missing
$\chrom \ge 6$ object simultaneously below by criticality (Kostochka-Yancey),
above by planarity-of-incidence ($K_{2,3}$-freeness), and at the abstract floor
by Folkman, to conclude that the object cannot be a small dense host.

\section{Discussion and the target object}
\label{sec:discussion}

The two preceding lines of argument pincer the missing ingredient from opposite
sides.
\begin{itemize}
  \item It cannot live in the realizable lineage
        (Sections~\ref{sec:census}-\ref{sec:criticality}): the lineage is
        vertex-critical and hence forcing-sterile, by the Essential-Pair Lemma.
  \item It cannot be a small dense Folkman-type host
        (Section~\ref{sec:codegree}): such hosts violate the codegree bound that
        every planar UDG must satisfy.
\end{itemize}
What remains is a single sharpened target.

\begin{quote}
\textbf{Target object.} A $6$-chromatic graph in the \emph{UDG-necessary class}
(simultaneously $K_4$-free and $K_{2,3}$-free), of order at least roughly $26$,
together with an explicit unit-distance embedding in $\R^2$.
\end{quote}

In-class membership ($K_4$-free and $K_{2,3}$-free) is a \emph{necessary}
condition for UDG-realizability, not a sufficient one. A graph passing both local
obstructions still requires an actual coordinate embedding to move
$\chrom(\R^2)$. The codegree ceiling, in particular, is a counting bound: the
true maximum number of unit distances among $n$ points sits well below it at small
$n$, so a graph that respects~\eqref{eq:codeg-ceiling} need not be realizable.

\paragraph{Honest limitations.} We list the caveats explicitly.
\begin{enumerate}
  \item \emph{The Essential-Pair Lemma is elementary and known.} Its corollary is
        Theorem~3.17 of Martin~\cite{Martin2009} (phrased there via
        ``implicit-edges'' and ``implicit-identities''), and it is implicit in
        standard critical-graph theory~\cite{JensenToft1995}. We include a proof
        only for self-containedness and make no priority claim.
  \item \emph{The $K_{2,3}$-free fact is classical.} It is the basis of the
        Spencer-Szemer\'edi-Trotter unit-distance edge bound. The novelty of
        Section~\ref{sec:codegree} is the synthesis with Kostochka-Yancey and the
        Folkman number, not the ingredients.
  \item \emph{The density observation is empirical.} The claim that reachable
        $K_4$-free $6$-critical hosts sit at $m/n \approx 4.2$ is an observation
        over roughly $35$ cores from one construction family, not a theorem about
        all $K_4$-free $6$-critical graphs. Only the $n \le 12$ exclusion
        (Proposition~\ref{prop:core13}) and the Folkman-floor exclusion are
        rigorous.
  \item \emph{No bound on $\chrom(\R^2)$ moved.} This note is negative and
        explanatory. It tightens the description of where the missing object can
        live; it does not produce the object.
\end{enumerate}

\paragraph{Outlook.} The results suggest a clear procedural lesson: generate
\emph{inside} the UDG-necessary class rather than above it. Manufacturing
$K_4$-free $6$-chromatic graphs and cutting them down fails the codegree bound;
the alternative is to enforce $K_4$-freeness and $K_{2,3}$-freeness as hard
invariants from the start, seeding from a real $\chrom = 5$ UDG (which lies in the
class automatically) and adding only codegree-safe and $K_4$-safe structure until
the chromatic number is forced upward. Whether that class forces $\chrom \ge 6$ at
some attainable order, or instead caps at $\chrom = 5$ (the UDG analogue of
$\chrom(\Q^2) = 2$), is the open question this framing isolates.

Preliminary computational probes of this question are suggestive but
inconclusive, and they approach the target object from three independent
directions, none of which reaches it. \emph{(i) Growth from below.} A greedy
search that grows a graph from the empty graph by adding only codegree-safe and
$K_4$-safe edges caps at $\chrom = 4$ across orders $17 \le n \le 64$: it never
even reaches $\chrom = 5$. \emph{(ii) Annealing.} A local search with both edge
additions and deletions (simulated annealing over graphs kept $K_4$-free and
$K_{2,3}$-free, minimizing the number of proper $5$-colorings) over orders
$26 \le n \le 40$, with on the order of $20{,}000$ restarts, reaches in-class
$\chrom = 5$ at orders as small as $n = 29$ but never $\chrom = 6$, with the edge
density plateauing near $84\%$ of the codegree
ceiling~\eqref{eq:codeg-ceiling} at every $n$. \emph{(iii) Repair from above.} The
top-down probe of Section~\ref{sec:codegree}, which starts from a genuine
$\chrom = 6$ graph ($M^3(C_5)$) and repairs its codegree violations while holding
$\chrom = 6$, stalls with the violations intact in the $6$-critical core. All
three are heuristic, not non-existence results: a $\chrom \ge 6$ in-class graph may
live at larger $n$, or at configurations none of the three local searches reach.
But together they triangulate the same conclusion---the target object, if it
exists in this range, is not easily constructed by local moves from any natural
seed.

\section*{Acknowledgments}
This is a note from an ongoing computational program on the Hadwiger-Nelson
problem. The realizable lineage graphs studied here derive from the
de~Grey~\cite{deGrey2018} and Polymath16~\cite{Polymath16} constructions. The
exhaustive forcing census and the criticality scans were carried out on a single
workstation using standard complete SAT solvers (CaDiCaL, MapleChrono, and
Glucose, invoked through the PySAT toolkit) together with a purpose-built
symmetry-breaking colorability solver; the local searches of
Section~\ref{sec:discussion} were run under PyPy. No external funding supported
this work.

\bibliographystyle{plain}
\bibliography{refs}

\end{document}
